%%%%%%%%%%%%%%%%%%%%%%%%%%%%%%%%%%%%%%%%%%%%%%%%%%%
% CHƯƠNG 6: THỬ NGHIỆM THỰC TẾ VÀ ĐÁNH GIÁ KẾT QUẢ
%%%%%%%%%%%%%%%%%%%%%%%%%%%%%%%%%%%%%%%%%%%%%%%%%%%
\section*{CHƯƠNG 6. THỬ NGHIỆM THỰC TẾ VÀ ĐÁNH GIÁ KẾT QUẢ}
\addcontentsline{toc}{section}{\numberline{}CHƯƠNG 6. THỬ NGHIỆM THỰC TẾ VÀ ĐÁNH GIÁ KẾT QUẢ}
\setcounter{section}{6}
\setcounter{subsection}{0}
\setcounter{figure}{0}
\setcounter{table}{0}

\subsection{Danh mục linh kiện thiết kế chế tạo (BOM)}
Danh sách toàn bộ các linh kiện chính được sử dụng để thiết kế và chế tạo sản phẩm bo mạch đeo tay được tổng hợp chi tiết trong Bảng \ref{tab:bom_linhkien_detail}.

\begin{table}[H]
    \centering
    \caption{Danh mục linh kiện chính sử dụng trong thiết kế bo mạch}
    \label{tab:bom_linhkien_detail}
    \fontsize{11pt}{14pt}\selectfont
    \renewcommand{\baselinestretch}{1.15}
    \renewcommand{\arraystretch}{1.15}
    \begin{tabularx}{\textwidth}{|c|X|X|c|X|}
        \hline
        \textbf{STT} & \textbf{Tên linh kiện} & \textbf{Mã hiệu / Thông số} & \textbf{SL} & \textbf{Chức năng trong mạch} \\ \hline
        1 & Vi điều khiển SoC & ESP32-C3-WROOM-02 & 1 & Khối xử lý trung tâm, lọc số DSP và truyền Wi-Fi \\ \hline
        2 & Cảm biến nhịp tim & MAX30102 & 1 & Thu nhận tín hiệu quang học PPG thể tích máu động mạch \\ \hline
        3 & Cảm biến nhiệt ẩm & HDC2022 & 1 & Đo nhiệt độ da và độ ẩm cơ thể người đeo \\ \hline
        4 & Màn hình hiển thị & OLED SSD1306 1.3" & 1 & Hiển thị thông số tại chỗ và trạng thái kết nối mạng \\ \hline
        5 & IC quản lý sạc nguồn & TP4056 & 1 & Quản lý sạc pin Lithium 3.7V an toàn qua USB \\ \hline
        6 & IC bảo vệ pin sạc & DW01A & 1 & Giám sát bảo vệ pin chống quá sạc, quá xả, quá dòng \\ \hline
        7 & Mosfet bảo vệ nguồn & FS8205A & 1 & Đóng cắt mạch xả sạc bảo vệ theo lệnh của DW01A \\ \hline
    \end{tabularx}
\end{table}

\subsection{Quá trình thử nghiệm thực tế}
Để đánh giá độ tin cậy thực tế của hệ thống, thiết bị phần cứng đeo tay được thử nghiệm đo đạc thực tế trên cơ thể 5 đối tượng tình nguyện khỏe mạnh có độ tuổi từ 20 đến 25 tuổi.
\begin{enumerate}
    \item \textbf{Cấu hình mạng ban đầu:} Thiết bị chưa có cấu hình mạng Wi-Fi tự động phát SSID \texttt{ESP32\_Setup}. Người dùng kết nối điện thoại vào và nhập Wi-Fi của phòng Lab. Thiết bị lưu Preferences thành công, reboot và kết nối mạng Wi-Fi thành công chỉ trong vòng 8 giây.
    \item \textbf{Thu thập dữ liệu tại chỗ:} Khi người dùng đeo đồng hồ vào cổ tay, cảm biến MAX30102 phát hiện áp lực tiếp xúc, đèn LED đỏ và hồng ngoại bật sáng bắt đầu thu thập sóng PPG. Màn hình OLED hiển thị chính xác giờ hệ thống đồng bộ NTP, giá trị nhịp tim, $SpO_2$, nhiệt độ da và độ ẩm.
    \item \textbf{Đồng bộ hóa dữ liệu lên server:} Dữ liệu sinh hiệu được đóng gói định dạng JSON gửi đều đặn mỗi 2 giây lên máy chủ API. Dashboard Web nhận dữ liệu và cập nhật biểu đồ lịch sử trơn tru mà không có hiện tượng mất gói tin hay nghẽn kết nối.
\end{enumerate}

\subsection{Đối sánh kết quả đo thực nghiệm và Đánh giá sai số}
Để đánh giá độ chính xác khoa học của thuật toán bù sai số cổ tay đã phát triển, tôi thực hiện đo đối sánh song song các chỉ số nhịp tim và $SpO_2$ giữa thiết bị đeo tay tự chế tạo (đeo ở cổ tay) với thiết bị đo y tế chuyên dụng Beurer PO30 (kẹp ngón tay) trên cùng một đối tượng đo tại 10 thời điểm đo khác nhau ở trạng thái nghỉ ngơi tĩnh và sau khi vận động nhẹ. Kết quả thực nghiệm đối sánh được ghi nhận chi tiết ở Bảng \ref{tab:comparison_data}.

\begin{table}[H]
    \centering
    \caption{Bảng đối sánh dữ liệu đo giữa thiết bị đeo tay và Beurer PO30}
    \label{tab:comparison_data}
    \fontsize{11pt}{14pt}\selectfont
    \renewcommand{\baselinestretch}{1.15}
    \renewcommand{\arraystretch}{1.15}
    \begin{tabularx}{\textwidth}{|c|X|X|X|X|p{3.5cm}|}
        \hline
        \textbf{Lần đo} & \textbf{Nhịp tim thiết bị (BPM)} & \textbf{Nhịp tim Beurer PO30} & \textbf{$SpO_2$ thiết bị (\%)} & \textbf{$SpO_2$ Beurer PO30} & \textbf{Trạng thái đối tượng} \\ \hline
        1 & 72 & 71 & 98.2 & 98.0 & Nghỉ ngơi tĩnh \\ \hline
        2 & 74 & 75 & 97.6 & 98.0 & Nghỉ ngơi tĩnh \\ \hline
        3 & 68 & 67 & 98.5 & 99.0 & Nghỉ ngơi tĩnh \\ \hline
        4 & 70 & 71 & 97.8 & 98.0 & Nghỉ ngơi tĩnh \\ \hline
        5 & 75 & 73 & 97.4 & 97.0 & Nghỉ ngơi tĩnh \\ \hline
        6 & 98 & 96 & 96.8 & 97.0 & Sau vận động nhẹ \\ \hline
        7 & 105 & 102 & 95.8 & 96.0 & Sau vận động nhẹ \\ \hline
        8 & 112 & 109 & 95.4 & 96.0 & Sau vận động nhẹ \\ \hline
        9 & 88 & 86 & 97.2 & 97.0 & Sau vận động nhẹ \\ \hline
        10 & 80 & 79 & 97.5 & 98.0 & Sau vận động nhẹ \\ \hline
    \end{tabularx}
\end{table}

Để đánh giá sai số tuyệt đối trung bình (MAE -- Mean Absolute Error) giữa hai thiết bị, tôi áp dụng công thức toán học:
\begin{equation}
    MAE = \frac{1}{N} \sum_{i=1}^{N} |X_{\text{device},i} - X_{\text{ref},i}|
\end{equation}
Kết quả tính toán sai số tuyệt đối trung bình từ dữ liệu thực nghiệm:
\begin{itemize}
    \item \textbf{Sai số nhịp tim trung bình:}
    \begin{equation}
        MAE_{\text{HeartRate}} = \frac{1 + 1 + 1 + 1 + 2 + 2 + 3 + 3 + 2 + 1}{10} = 1.7\text{ BPM}
    \end{equation}
    \item \textbf{Sai số $SpO_2$ trung bình:}
    \begin{equation}
        MAE_{SpO_2} = \frac{0.2 + 0.4 + 0.5 + 0.2 + 0.4 + 0.2 + 0.2 + 0.6 + 0.2 + 0.5}{10} \approx 0.34\%
    \end{equation}
\end{itemize}

Sai số nhịp tim dưới 2 BPM và sai số $SpO_2$ dưới 1\% hoàn toàn đạt các tiêu chuẩn kỹ thuật cho phép đối với các thiết bị giám sát sức khỏe cá nhân phi lâm sàng, chứng minh tính hiệu quả vượt trội của thuật toán bù sai số thích ứng đo ở cổ tay đã thiết kế.

Thời lượng hoạt động thực tế của pin sạc Lithium 3.7V dung lượng 500mAh đo được đạt mức khoảng 18.5 giờ hoạt động liên tục ở chế độ truyền dữ liệu qua Wi-Fi 2 giây một lần. Hệ thống cảnh báo trên Web Dashboard cũng tự động kích hoạt và hiển thị cảnh báo đỏ trực quan tức thời khi phát hiện sinh hiệu bệnh nhân bất thường vượt quá các ngưỡng y sinh an toàn đã thiết lập.
